\documentclass[12pt,letterpaper]{article}

% -------------------------------------------------
% Formato general tipo APA 7.ª edición
% -------------------------------------------------
\usepackage[utf8]{inputenc}
\usepackage[T1]{fontenc}
\usepackage[spanish,es-tabla]{babel}

% Configuración para usar Helvetica (clon directo de Arial) en TODO el documento
\usepackage{helvet}
\renewcommand{\familydefault}{\sfdefault}
\shorthandoff{"}

\usepackage{geometry}
\usepackage{setspace}
\usepackage{indentfirst}
\usepackage[table]{xcolor}
\usepackage{hyperref}
\usepackage{xurl}
\usepackage{longtable}
\usepackage{array}
\usepackage{booktabs}
\usepackage{caption}
\usepackage{needspace}
\usepackage{enumitem}
\usepackage{fancyhdr}
\usepackage{titlesec}
\usepackage{listings}
\usepackage{graphicx}
\usepackage{ragged2e}
\usepackage{float}
\usepackage{etoolbox}
\usepackage{tocloft}

% -------------------------------------------------
% Mejora visual de tablas y figuras
% -------------------------------------------------
\definecolor{tableStripe}{HTML}{F2F4F7}
\renewcommand{\arraystretch}{1.20}
\setlength{\arrayrulewidth}{0.8pt}
\setlength{\tabcolsep}{8pt}
\setlength{\LTpre}{0.5em}
\setlength{\LTpost}{0.5em}

% Se fuerza a que el título autogenerado de tablas y figuras quede ARRIBA
\captionsetup[table]{font={small,sf},labelfont=bf,position=above,skip=6pt}
\captionsetup[figure]{font={small,sf},labelfont=bf,position=above,skip=6pt}

\AtBeginEnvironment{longtable}{%
  \setlength{\tabcolsep}{5pt}%
  \rowcolors{2}{tableStripe}{white}%
}

% Márgenes APA
\geometry{margin=1in}
\setlength{\headheight}{15pt}

% Interlineado
\doublespacing

% Sangría
\setlength{\parindent}{0.5in}
\setlength{\parskip}{0pt}

% Formateo estricto de listas bajo estándares tipográficos
\setlist[itemize]{leftmargin=0.5in, topsep=0pt, partopsep=0pt, parsep=0pt, itemsep=4pt}
\setlist[enumerate]{leftmargin=0.5in, topsep=0pt, partopsep=0pt, parsep=0pt, itemsep=4pt}

% Alineación APA
\raggedright

% -------------------------------------------------
% Encabezado y numeración
% -------------------------------------------------
\pagestyle{fancy}
\fancyhf{}
\rhead{\thepage}
\renewcommand{\headrulewidth}{0pt}

% -------------------------------------------------
% Colores y enlaces
% -------------------------------------------------
\definecolor{apaBlue}{HTML}{1F4E79}

\hypersetup{
    colorlinks=true,
    linkcolor=black,
    urlcolor=apaBlue,
    citecolor=black,
    pdftitle={Manual de Código: Solicitudes Académicas FESC},
    pdfauthor={Desarrolladores FESC}
}

% -------------------------------------------------
% Jerarquía de títulos (Todos en tamaño 12pt / \normalsize)
% -------------------------------------------------
\titleformat{\section}
  {\normalfont\bfseries\normalsize}
  {\thesection.}{0.5em}{}

\titleformat{\subsection}
  {\normalfont\bfseries\normalsize}
  {\thesubsection}{0.5em}{}

\titleformat{\subsubsection}
  {\normalfont\bfseries\itshape\normalsize}
  {\thesubsubsection}{0.5em}{}

\titlespacing*{\section}{0pt}{1.5em}{1em}
\titlespacing*{\subsection}{0pt}{1.2em}{0.8em}
\titlespacing*{\subsubsection}{0pt}{1em}{0.6em}

% Cada capítulo inicia en hoja nueva de manera limpia
\pretocmd{\section}{\clearpage}{}{}

% -------------------------------------------------
% Código fuente (SIN CUADROS, SIN FONDO, FUENTE ARIAL/HELVETICA TAMAÑO 12)
% -------------------------------------------------
\lstdefinelanguage{typescript}{
    keywords={const,let,var,function,return,export,import,from,type,interface,extends,if,else,await,async,null,true,false,default,new,class},
    sensitive=true,
    comment=[l]{//},
    morecomment=[s]{/*}{*/},
    morestring=[b]',
    morestring=[b]",
    morestring=[b]`
}

\lstdefinestyle{codeblock}{
    backgroundcolor=\color{white},
    frame=none,
    basicstyle=\normalfont\normalsize,
    breaklines=true,
    breakatwhitespace=false,
    columns=fullflexible,
    keepspaces=true,
    showstringspaces=false,
    literate={á}{{\'a}}1 {é}{{\'e}}1 {í}{{\'i}}1 {ó}{{\'o}}1 {ú}{{\'u}}1
             {Á}{{\'A}}1 {É}{{\'E}}1 {Í}{{\'I}}1 {Ó}{{\'O}}1 {Ú}{{\'U}}1
             {ñ}{{\~n}}1 {Ñ}{{\~N}}1 {ü}{{\"u}}1 {Ü}{{\"U}}1
             {→}{{$\to$}}1 {≥}{{$\geq$}}1 {≤}{{$\leq$}}1,
    tabsize=2
}
\lstset{style=codeblock}

% Comando seguro para imprimir nombres técnicos con guiones bajos sin romper el entorno matemático
\newcommand{\code}[1]{\texorpdfstring{\begingroup\normalfont\normalsize\def\_{\textunderscore\allowbreak}#1\endgroup}{#1}}

% -------------------------------------------------
% Figuras Corregidas: Nombre arriba (vía caption), imagen, descripción abajo
% -------------------------------------------------
\newcommand{\diagramfigurehere}[4]{%
\begin{figure}[H]
\centering
\caption{#2}
\label{#4}
\vspace{0.2cm}

\IfFileExists{#1}{%
    \includegraphics[width=0.9\textwidth,height=0.65\textheight,keepaspectratio]{#1}%
}{%
    \fbox{%
        \begin{minipage}[c][2.5in][c]{0.88\textwidth}
        \centering
        \textbf{Espacio reservado para el diagrama}\\[0.5em]
        Inserta aquí la imagen correspondiente en la ruta:\\[0.3em]
        \texttt{#1}
        \end{minipage}%
    }%
}

\vspace{0.2cm}
{\small \textit{Descripción de la figura: #3}}
\end{figure}
}

\newcommand{\diagramfigurefullpage}[4]{%
\begin{figure}[H]
\centering
\caption{#2}
\label{#4}
\vspace{0.2cm}

\IfFileExists{#1}{%
    \includegraphics[width=\textwidth,height=0.82\textheight,keepaspectratio]{#1}%
}{%
    \fbox{%
        \begin{minipage}[c][4in][c]{0.95\textwidth}
        \centering
        \textbf{Espacio reservado para el diagrama}\\[0.5em]
        Inserta aquí la imagen correspondiente en la ruta:\\[0.3em]
        \texttt{#1}
        \end{minipage}%
    }%
}

\vspace{0.2cm}
{\small \textit{Descripción de la figura: #3}}
\end{figure}
}

\begin{document}

% ==================================================================
% PORTADA ACTUALIZADA
% ==================================================================
\begin{titlepage}
    \centering
    \vspace*{0.5in}

    % 1. Nombres completos de los estudiantes actualizados
    {\normalsize\bfseries
    Laura Valentina Rozzo Espinel\\
    Nelly Fabiola Cano Oviedo\\
    Nestor Ivan Granados Valenzuela\\
    Nick Alejandro Ortega Mendez\\
    Santiago Alejandro Medina Ortega\\
    Valery Gabriela Alarcon Peña
    \par}
    \vspace{0.6in}

    % 2. Programa académico
    {\normalsize Ingeniería de Software \par}
    \vspace{0.4in}

    % 3. Nombre completo de la FESC
    {\normalsize Fundación de Estudios Superiores Comfanorte (FESC) \par}
    \vspace{0.6in}

    % 4. Nombre de la asignatura
    {\normalsize\bfseries Programación Orientada a Objetos \par}
    \vspace{0.4in}

    % 5. Identificación del Proyecto (PPA) y Título del Documento
    {\normalsize Proyecto Pedagógico de Aula (PPA) \par}
    \vspace{0.5in}
    {\normalsize\bfseries MANUAL DE CÓDIGO:\ SISTEMA DE SOLICITUDES ACADÉMICAS FESC\par}
    \vspace{0.6in}

    % 6. Nombre completo del docente actualizado
    {\normalsize Ing. Angel Ureña \par}

    \vfill
    {\normalsize San José de Cúcuta, Mayo de 2026 \par}
\end{titlepage}

\newpage
\tableofcontents

\newpage
\addcontentsline{toc}{section}{Índice de tablas}
\listoftables

\newpage
\addcontentsline{toc}{section}{Índice de figuras}
\listoffigures

% =================================================
\section*{Introducción}
\addcontentsline{toc}{section}{Introducción}
% =================================================

El presente manual de código documenta la arquitectura, los componentes principales, las convenciones y el flujo de trabajo del proyecto \textbf{Sistema de Solicitudes Académicas FESC}. El sistema es una aplicación web desarrollada en Java EE con servlets, JSP, MySQL y Apache Tomcat, orientada a la gestión de solicitudes académicas de estudiantes y la administración de trámites internos.

El documento está dirigido a desarrolladores e integradores del proyecto, con la finalidad de facilitar el mantenimiento, la extensión y el despliegue correcto de la solución. Incluye la descripción de la estructura de carpetas, rutas, configuración local, lógica principal, estándares de código, persistencia y los diagramas necesarios para visualización arquitectónica.

% =================================================
\section{Visión General}
% =================================================

Este sistema gestiona el ciclo completo de solicitudes académicas de los estudiantes, desde el registro y autenticación, hasta la aprobación, seguimiento y auditoría de retrasos por parte de los administradores. La aplicación separa claramente las áreas de estudiante y administrador, y utiliza un backend Java que se comunica con una base de datos MySQL a través de JDBC.

\Needspace{6\baselineskip}
\begin{longtable}{|>{\bfseries}p{0.25\textwidth}|p{0.65\textwidth}|}
\caption{Áreas principales del sistema}\\
\hline
Área & Descripción \\
\hline
Portal de estudiante & Permite a los estudiantes autenticarse, ver su perfil, crear solicitudes, consultar su estado y recibir mensajes de seguimiento. \\
\hline
Panel administrativo & Permite a los administradores gestionar estudiantes, programas, solicitudes, tipos de solicitud, responsables asignados y generar reportes. \\
\hline
Infraestructura de datos & Base de datos MySQL que almacena sedes, jornadas, programas académicos, estudiantes, administradores, solicitudes, mensajes y auditorías. \\
\hline
\end{longtable}

La aplicación se ejecuta en un contenedor Java EE compatible (Apache Tomcat) y utiliza JSP para renderizar vistas, servlets para manejar la lógica HTTP y un DAO centralizado para el acceso a datos.

\subsection*{Flujo de datos de alto nivel}

\diagramfigurehere
{diagramas/flujo-datos-alto-nivel.png}
{Flujo de datos de alto nivel del sistema}
{Relación entre el navegador, los servlets Java, el conjunto de JSP, la capa DAO y la base de datos MySQL.}
{fig:flujo-datos-alto-nivel}

% =================================================
\section{Stack Tecnológico}
% =================================================

Esta sección describe las tecnologías y herramientas usadas en el proyecto real.

\subsection{Tecnologías principales}

\Needspace{8\baselineskip}
\begin{longtable}{|p{0.32\textwidth}|p{0.18\textwidth}|p{0.40\textwidth}|}
\caption{Tecnologías principales del proyecto}\\
\hline
\textbf{Tecnología} & \textbf{Versión / Tipo} & \textbf{Uso} \\
\hline
Java SE & 8+ & Lenguaje de programación backend y lógica de negocio. \\
\hline
Java Servlet & 3.1 & Gestión de peticiones HTTP y controladores. \\
\hline
JSP & 2.x & Vistas dinámicas del frontend. \\
\hline
MySQL / MariaDB & 8.x recomendado & Motor de base de datos relacional. \\
\hline
JDBC & MySQL Connector/J & Conexión entre la aplicación Java y la base de datos. \\
\hline
Apache Tomcat & 8.x / 9.x & Contenedor de servlets para despliegue. \\
\hline
Apache Ant & 1.10+ & Herramienta de construcción y empaquetado. \\
\hline
NetBeans & IDE compatible & Desarrollo y ejecución del proyecto Java EE. \\
\hline
\end{longtable}

\subsection{Herramientas de desarrollo}

\Needspace{6\baselineskip}
\begin{longtable}{|p{0.32\textwidth}|p{0.18\textwidth}|p{0.40\textwidth}|}
\caption{Herramientas de desarrollo utilizadas}\\
\hline
\textbf{Herramienta} & \textbf{Versión / Uso} & \textbf{Descripción} \\
\hline
Apache Ant & Build tool & Compila el código, copia recursos y genera el despliegue. \\
\hline
NetBeans & IDE Java EE & Entorno de desarrollo y ejecución con soporte Tomcat. \\
\hline
MySQL Workbench / cliente SQL & Administración de BD & Gestión de esquema y datos. \\
\hline
VS Code & Editor de código & Edición de archivos y configuración opcional de tareas Ant. \\
\hline
\end{longtable}

% =================================================
\section{Arquitectura del Proyecto}
% =================================================

La arquitectura del proyecto está organizada en capas ligeras: presentación JSP, controladores servlet, acceso a datos DAO, modelos de datos y utilidades. El frontend se compone de archivos JSP y recursos estáticos bajo la carpeta \code{web/}, mientras que el backend Java reside en \code{src/java/}.

\subsection{Estructura de carpetas}

\Needspace{10\baselineskip}
\begin{lstlisting}[language={}] 
solicitudes-academicas/
|---- build.xml
|---- README.md
|---- docs/
|    |---- INSTRUCCIONES_CONFIGURACION.md
|---- nbproject/
|---- sql/
|    `---- database.sql
|---- src/
|    `---- java/
|        |---- controller/
|        |---- dao/
|        |---- filter/
|        |---- model/
|        |---- util/
|        |---- database.properties.example
|        `---- database.properties
|---- web/
|    |---- WEB-INF/
|    |    |---- web.xml
|    |    `---- classes/
|    |---- admin/
|    |---- student/
|    |---- assets/
`---- tmp-classes/
\end{lstlisting}

\subsection{Organización interna de paquetes}

\Needspace{6\baselineskip}
\begin{lstlisting}[language={}]
src/java/
|---- controller/    # Servlets que reciben peticiones y redirigen a vistas
|---- dao/           # Acceso a datos y operaciones contra MySQL
|---- model/         # Objetos de dominio que representan tablas
|---- util/          # Utilidades compartidas como hashing y conexion DB
|---- filter/        # Filtros HTTP para comportamiento transversal
\end{lstlisting}

\subsection{Diagrama de módulos y dependencias}

\diagramfigurehere
{diagramas/modulos-y-dependencias.png}
{Diagrama de módulos y dependencias del proyecto}
{Relación entre los componentes de presentación, controladores, DAO, modelos y utilidades Java.}
{fig:modulos-dependencias}

\subsection{Sistema de rutas}

Estas son las rutas principales definidas en los servlets del proyecto.

\Needspace{12\baselineskip}
\begin{longtable}{|p{0.32\textwidth}|p{0.38\textwidth}|p{0.18\textwidth}|}
\caption{Rutas principales del sistema}\\
\hline
\textbf{Ruta} & \textbf{Servlet} & \textbf{Protegida} \\
\hline
\code{/login} & \code{LoginServlet} & No \\
\hline
\code{/register} & \code{RegisterServlet} & No \\
\hline
\code{/forgot-password} & \code{ForgotPasswordServlet} & No \\
\hline
\code{/recover-password} & \code{RecoverPasswordServlet} & No \\
\hline
\code{/logout} & \code{LogoutServlet} & Sí \\
\hline
\code{/student/dashboard} & \code{StudentDashboardServlet} & Sí \\
\hline
\code{/student/profile} & \code{StudentProfileServlet} & Sí \\
\hline
\code{/student/requests} & \code{MyRequestsServlet} & Sí \\
\hline
\code{/student/request-detail} & \code{RequestDetailServlet} & Sí \\
\hline
\code{/student/new-request} & \code{NewRequestServlet} & Sí \\
\hline
\code{/student/edit-request} & \code{EditRequestServlet} & Sí \\
\hline
\code{/student/cancel-request} & \code{CancelRequestServlet} & Sí \\
\hline
\code{/admin/dashboard} & \code{AdminDashboardServlet} & Sí \\
\hline
\code{/admin/profile} & \code{AdminProfileServlet} & Sí \\
\hline
\code{/admin/admin-users} & \code{AdminUsersServlet} & Sí \\
\hline
\code{/admin/students} & \code{AdminStudentsServlet} & Sí \\
\hline
\code{/admin/programs} & \code{AdminProgramsServlet} & Sí \\
\hline
\code{/admin/request-types} & \code{AdminRequestTypesServlet} & Sí \\
\hline
\code{/admin/requests} & \code{AdminRequestsServlet} & Sí \\
\hline
\code{/admin/request-detail} & \code{AdminRequestDetailServlet} & Sí \\
\hline
\code{/admin/reports} & \code{AdminReportsServlet} & Sí \\
\hline
\end{longtable}

% =================================================
\section{Guía de Configuración}
% =================================================

Esta sección describe los requisitos y pasos para instalar, configurar y ejecutar el sistema en un entorno local de desarrollo.

\subsection{Prerrequisitos}

Antes de ejecutar el sistema, el equipo debe contar con:

\begin{itemize}
    \item Java JDK 8 o superior.
    \item Apache Tomcat instalado localmente.
    \item MySQL o MariaDB con la base de datos \code{fesc\_solicitudes}.
    \item Apache Ant para compilar y empaquetar el proyecto.
\end{itemize}

\subsection{Instalación paso a paso}

\Needspace{10\baselineskip}
\begin{lstlisting}[language=bash]
# 1. Clonar el repositorio
git clone <url-del-repo>
cd solicitudes-academicas

# 2. Crear archivo de configuracion local
cp src/java/database.properties.example src/java/database.properties

# 3. Ajustar credenciales MySQL
# Editar src/java/database.properties

# 4. Compilar con Ant
ant clean
ant

# 5. Desplegar en Tomcat
# Copiar WAR o carpeta compilada a TOMCAT_HOME/webapps/
\end{lstlisting}

\subsection{Variables de entorno y configuración de DB}

El archivo \code{src/java/database.properties} contiene los parámetros de conexión a MySQL:

\Needspace{4\baselineskip}
\begin{lstlisting}[language={}]
db.url=jdbc:mysql://localhost:3306/fesc_solicitudes?useSSL=false&allowPublicKeyRetrieval=true&serverTimezone=America/Bogota
db.user=root
db.password=TU_CONTRASENA
\end{lstlisting}

Este archivo no debe subirse al repositorio; el archivo \code{database.properties.example} sirve como plantilla.

\subsection{Scripts disponibles}

\Needspace{5\baselineskip}
\begin{longtable}{|p{0.22\textwidth}|p{0.28\textwidth}|p{0.40\textwidth}|}
\caption{Scripts disponibles en el proyecto}\\
\hline
\textbf{Script} & \textbf{Comando} & \textbf{Descripción} \\
\hline
Clean & \code{ant clean} & Elimina artefactos de compilación previos. \\
\hline
Build & \code{ant} & Compila el proyecto y genera la salida de despliegue. \\
\hline
\end{longtable}

% =================================================
\section{Documentación de Lógica Principal}
% =================================================

Esta sección documenta los archivos y procesos más relevantes del código Java.

\subsection{\texorpdfstring{\code{src/java/dao/DBConnection.java}}{src/java/dao/DBConnection.java}}

Este archivo centraliza la conexión JDBC con MySQL. Lee \code{database.properties} de la ruta de clases y carga el driver \code{com.mysql.cj.jdbc.Driver}.

\Needspace{4\baselineskip}
\begin{lstlisting}[language=java]
public static Connection getConnection() throws SQLException {
    return DriverManager.getConnection(URL, USER, PASSWORD);
}
\end{lstlisting}

La regla principal es mantener una sola fuente de verdad para la conexión y no instanciar drivers en múltiples clases.

\subsection{\texorpdfstring{\code{src/java/util/PasswordHasher.java}}{src/java/util/PasswordHasher.java}}

Este helper gestiona el hashing seguro de contraseñas usando SHA-256 y una sal Base64.

\Needspace{4\baselineskip}
\begin{lstlisting}[language=java]
public static String secureHash(String password) {
    String salt = generateSalt();
    String hash = hashPassword(password, salt);
    return salt + ":" + hash;
}
\end{lstlisting}

\subsection{Acceso a datos: DAOs principales}

Los DAOs contienen la lógica de lectura y escritura contra MySQL. Entre los más importantes están:

\begin{itemize}
    \item \code{AdminDAO.java}
    \item \code{StudentDAO.java}
    \item \code{SolicitudMensajeDAO.java}
    \item \code{TipoSolicitudDAO.java}
    \item \code{ProgramaAcademicoDAO.java}
\end{itemize}

Estos archivos encapsulan consultas SQL y mapeos entre resultados y objetos de dominio.

\subsection{Servicio de solicitudes}

Ubicación principal: \code{src/java/controller/} y \code{src/java/dao/}.

El flujo principal de solicitudes académicas sigue estos pasos:

\begin{enumerate}
    \item El estudiante envía un formulario a \code{NewRequestServlet}.
    \item El servlet valida campos y crea el registro en \code{solicitud} mediante \code{StudentDAO}.
    \item El administrador revisa la solicitud en \code{AdminRequestsServlet}.
    \item Las respuestas y el responsable se actualizan en \code{solicitud} mediante \code{AdminRequestDetailServlet}.
    \item El historial se registra en \code{solicitud\_mensaje} y \code{auditoria\_retraso}.
\end{enumerate}

\subsection*{Flujo de estados del formulario de solicitud}

\diagramfigurehere
{diagramas/estado-formulario-solicitud.png}
{Diagrama de estados del formulario de solicitud}
{Transición de estados del formulario de creación/edición de solicitud desde la carga inicial hasta el envío y manejo de errores.}
{fig:estado-formulario-solicitud}

\subsection{Protección de rutas y sesiones}

La aplicación utiliza filtros y sesiones HTTP para proteger el acceso a rutas administrativas y de estudiante. El servlet de logout cierra la sesión activa.

\subsubsection*{Flujo de autenticación}

\diagramfigurehere
{diagramas/flujo-autenticacion.png}
{Flujo general del módulo de autenticación}
{Proceso de inicio de sesión, validación de credenciales, creación de sesión y acceso a rutas protegidas.}
{fig:flujo-autenticacion}

% =================================================
\section{Estándares de Código}
% =================================================

Esta sección define las reglas de estilo y organización del código del proyecto actual.

\subsection{Convenciones de nombrado}

\Needspace{10\baselineskip}
\begin{longtable}{|p{0.30\textwidth}|p{0.30\textwidth}|p{0.30\textwidth}|}
\caption{Convenciones de nombrado del proyecto}\\
\hline
\textbf{Elemento} & \textbf{Convención} & \textbf{Ejemplo} \\
\hline
Clases Java & \code{PascalCase} & \code{StudentDAO} \\
\hline
Servlets & \code{PascalCase} & \code{LoginServlet} \\
\hline
Modelos & \code{PascalCase} & \code{ProgramaAcademico} \\
\hline
Métodos & \code{camelCase} & \code{getConnection} \\
\hline
Variables locales & \code{camelCase} & \code{studentId} \\
\hline
Constantes & \code{UPPER\_SNAKE\_CASE} & \code{DB\_URL} \\
\hline
Archivos JSP & \code{kebab-case.jsp} & \code{forgot\_password.jsp} \\
\hline
Archivos Java & \code{PascalCase.java} & \code{AdminReportsServlet.java} \\
\hline
Propiedades & \code{lowercase.with.dots} & \code{db.url} \\
\hline
\end{longtable}

\subsection{Estructura de un servlet}

Un servlet debe mantener un orden claro: inicialización, proceso de petición, delegación al DAO y despacho a JSP.

\Needspace{10\baselineskip}
\begin{lstlisting}[language=java]
@WebServlet("/login")
public class LoginServlet extends HttpServlet {
  protected void doPost(HttpServletRequest req, HttpServletResponse resp)
      throws ServletException, IOException {
    String email = req.getParameter("email");
    String password = req.getParameter("password");

    // Validar credenciales
    // Consultar StudentDAO / AdminDAO
    // Crear sesion
    // Redirigir a la pagina correspondiente
  }
}
\end{lstlisting}

\subsection{Manejo de excepciones}

Las excepciones SQL y de IO deben capturarse en el DAO o utilitario, y traducirse a mensajes controlados para el usuario o al log.

\subsection{Buenas prácticas de SQL}

\begin{itemize}
    \item Usar sentencias preparadas (\code{PreparedStatement}) en DAO para evitar inyección SQL.
    \item Mantener las consultas SQL en métodos específicos del DAO.
    \item Cerrar conexiones en bloques \code{try-with-resources}.
\end{itemize}

% =================================================
\section{Flujo de Trabajo}
% =================================================

Esta sección describe los procedimientos de desarrollo, extensión y publicación de cambios en el proyecto.

\subsection{Añadir una nueva funcionalidad de solicitud}

\Needspace{8\baselineskip}
\begin{lstlisting}[language={}]
1. Crear o actualizar la tabla SQL en sql/database.sql.
2. Generar el modelo Java en src/java/model/.
3. Anadir la logica de acceso en src/java/dao/.
4. Crear o extender el servlet en src/java/controller/.
5. Agregar las vistas JSP en web/ o web/student/ o web/admin/.
6. Probar la funcionalidad en Tomcat.
\end{lstlisting}

\subsection{Agregar una nueva ruta de servlet}

\begin{enumerate}
    \item Crear la clase Java en \code{src/java/controller/}.
    \item Anotar con \code{@WebServlet(name = "...", urlPatterns = {"/ruta"})}.
    \item Implementar \code{doGet} y/o \code{doPost}.
    \item Crear la vista JSP necesaria.
    \item Actualizar el menú o los enlaces de navegación.
\end{enumerate}

\subsection{Control de versiones}

\Needspace{5\baselineskip}
\begin{lstlisting}[language={}]
feature/alta-solicitud
fix/login-seguridad
chore/actualizar-configuracion-db
docs/manual-codigo
\end{lstlisting}

% =================================================
\section{Persistencia y Base de Datos}
% =================================================

El proyecto utiliza un esquema MySQL que define tablas maestras de sedes, jornadas y programas, además de tablas de usuarios, solicitudes, mensajes y auditoría.

\subsection{Modelo entidad-relación}

\diagramfigurefullpage
{diagramas/modelo-entidad-relacion.png}
{Modelo entidad-relación de la base de datos del sistema}
{Entidades principales del esquema MySQL, sus llaves primarias, llaves foráneas y relaciones de integridad.}
{fig:modelo-er}

\subsection{Esquema de tablas principales}

La base de datos incluye las siguientes tablas más importantes:

\begin{itemize}
    \item \code{sede}
    \item \code{jornada}
    \item \code{programa\_academico}
    \item \code{tipo\_solicitud}
    \item \code{estudiante}
    \item \code{administrador}
    \item \code{tipo\_solicitud\_responsable}
    \item \code{solicitud}
    \item \code{solicitud\_mensaje}
    \item \code{auditoria\_retraso}
\end{itemize}

\subsection{Relaciones clave}

\begin{itemize}
    \item \code{estudiante.programa\_id} $\rightarrow$ \code{programa\_academico.id}
    \item \code{estudiante.sede\_id} $\rightarrow$ \code{sede.id}
    \item \code{estudiante.jornada\_id} $\rightarrow$ \code{jornada.id}
    \item \code{solicitud.estudiante\_id} $\rightarrow$ \code{estudiante.id}
    \item \code{solicitud.tipo\_solicitud\_id} $\rightarrow$ \code{tipo\_solicitud.id}
    \item \code{solicitud.admin\_id} $\rightarrow$ \code{administrador.id}
    \item \code{solicitud.responsable\_id} $\rightarrow$ \code{administrador.id}
    \item \code{solicitud\_mensaje.solicitud\_id} $\rightarrow$ \code{solicitud.id}
    \item \code{auditoria\_retraso.solicitud\_id} $\rightarrow$ \code{solicitud.id}
    \item \code{auditoria\_retraso.admin\_id} $\rightarrow$ \code{administrador.id}
\end{itemize}

% =================================================
\section{Documentación Detallada de Controladores}
% =================================================

A continuación se documenta la responsabilidad principal de cada servlet (controlador) identificado en el proyecto, su ruta y las interacciones con DAOs y vistas.

\subsection*{LoginServlet}
\addcontentsline{toc}{subsection}{LoginServlet}
Ruta: \code{/login}. Archivo: \code{src/java/controller/LoginServlet.java}.
Responsabilidad: validar credenciales de administrativos y estudiantes (delegando a \code{AdminDAO} o \code{StudentDAO}), crear la sesión HTTP con los atributos de usuario y redirigir al dashboard correspondiente. Maneja mensajes de error por credenciales inválidas.

\subsection*{RegisterServlet}
\addcontentsline{toc}{subsection}{RegisterServlet}
Ruta: \code{/register}. Archivo: \code{src/java/controller/RegisterServlet.java}.
Responsabilidad: validar datos de registro (correo institucional, campos obligatorios), invocar \code{StudentDAO} para persistir el nuevo estudiante y redirigir a la pantalla de login o mostrar errores de validación.

\subsection*{ForgotPasswordServlet / RecoverPasswordServlet}
\addcontentsline{toc}{subsection}{ForgotPasswordServlet / RecoverPasswordServlet}
Rutas: \code{/forgot-password}, \code{/recover-password}. Archivos: \code{ForgotPasswordServlet.java}, \code{RecoverPasswordServlet.java}.
Responsabilidad: generar token de recuperación, almacenar token y expiración en la base de datos (vía \code{StudentDAO} o \code{AdminDAO}), y presentar la lógica para restablecer la contraseña. En la implementación de referencia el enlace se imprime en consola (simulador de email).

\subsection*{LogoutServlet}
\addcontentsline{toc}{subsection}{LogoutServlet}
Ruta: \code{/logout}. Archivo: \code{LogoutServlet.java}.
Responsabilidad: invalidar la sesión HTTP y redirigir a la página de login.

\subsection*{StudentDashboardServlet}
\addcontentsline{toc}{subsection}{StudentDashboardServlet}
Ruta: \code{/student/dashboard}. Archivo: \code{StudentDashboardServlet.java}.
Responsabilidad: recopilar estadísticas y la lista de solicitudes del estudiante (vía \code{StudentDAO} y \code{SolicitudDAO}), y despachar a la vista JSP del dashboard del estudiante.

\subsection*{StudentProfileServlet}
\addcontentsline{toc}{subsection}{StudentProfileServlet}
Ruta: \code{/student/profile}. Archivo: \code{StudentProfileServlet.java}.
Responsabilidad: mostrar y actualizar los datos del estudiante; delega operaciones de persistencia a \code{StudentDAO}.

\subsection*{MyRequestsServlet, NewRequestServlet, RequestDetailServlet, EditRequestServlet, CancelRequestServlet}
\addcontentsline{toc}{subsection}{MyRequests y gestión de solicitudes}
Rutas: \code{/student/requests}, \code{/student/new-request}, \code{/student/request-detail}, \code{/student/edit-request}, \code{/student/cancel-request}.
Responsabilidades:
\begin{itemize}
    \item \code{MyRequestsServlet}: listar las solicitudes del estudiante y aplicar filtros.
    \item \code{NewRequestServlet}: procesar la creación de una solicitud, validar campos y guardar adjuntos en \code{/uploads}; utiliza \code{StudentDAO} o \code{SolicitudDAO} para persistir datos.
    \item \code{RequestDetailServlet}: presentar detalle de solicitud, mensajes e historial; aceptar subida de archivos en mensajes si aplica.
    \item \code{EditRequestServlet}: permitir la edición de solicitudes en estados permitidos; aplica validaciones y persiste cambios.
    \item \code{CancelRequestServlet}: marcar la solicitud como anulada/cancelada y registrar auditoría.
\end{itemize}

\subsection*{AdminDashboardServlet, AdminRequestsServlet, AdminRequestDetailServlet}
\addcontentsline{toc}{subsection}{Admin: solicitudes y dashboard}
Rutas: \code{/admin/dashboard}, \code{/admin/requests}, \code{/admin/request-detail}.
Responsabilidades: mostrar métricas del sistema, listar solicitudes administrativas con filtros, gestionar el detalle de cada solicitud (cambiar estado, añadir respuesta, subir adjuntos desde admin). \code{AdminRequestDetailServlet} valida tipos de archivo y guarda adjuntos en \code{/uploads/messages}.

\subsection*{AdminUsersServlet, AdminStudentsServlet, AdminProgramsServlet, AdminRequestTypesServlet}
\addcontentsline{toc}{subsection}{Admin: gestión de recursos}
Rutas y responsabilidades:
\begin{itemize}
    \item \code{AdminUsersServlet}: gestionar cuentas administrativas (crear/editar/eliminar según rol).
    \item \code{AdminStudentsServlet}: administrar cuentas de estudiantes, buscar, editar y restablecer contraseñas.
    \item \code{AdminProgramsServlet}: CRUD de programas académicos.
    \item \code{AdminRequestTypesServlet}: CRUD y activación/desactivación de tipos de solicitud.
\end{itemize}

\subsection*{AdminReportsServlet}
\addcontentsline{toc}{subsection}{AdminReportsServlet}
Ruta: \code{/admin/reports}. Archivo: \code{AdminReportsServlet.java}.
Responsabilidad: construir vistas y datasets filtrados para reportes, y exportar resultados (CSV/PDF) según la vista.

% =================================================
\section{Modelos de Dominio}
% =================================================

Se listan los modelos principales que representan las tablas de la base de datos y su responsabilidad.

\subsection*{Student}
\addcontentsline{toc}{subsection}{Student}
Archivo: \code{src/java/model/Student.java}.
Responsabilidad: encapsular datos de un estudiante (id, nombre, apellido, email, password-hash, programaId, sedeId, jornadaId, etc.) y servir como DTO entre DAO y controladores.

\subsection*{Admin}
\addcontentsline{toc}{subsection}{Admin}
Archivo: \code{src/java/model/Admin.java}.
Responsabilidad: representar cuentas administrativas y sus atributos (id, nombre, email, rol), utilizado en autenticación y asignación de responsabilidades.

\subsection*{Solicitud}
\addcontentsline{toc}{subsection}{Solicitud}
Archivo: \code{src/java/model/Solicitud.java}.
Responsabilidad: representar una solicitud académica con sus campos principales (id, estudianteId, tipoSolicitudId, descripcion, estado, fechaCreacion, fechaActualizacion, adjuntos, responsable, adminAsignado).

\subsection*{SolicitudMensaje}
\addcontentsline{toc}{subsection}{SolicitudMensaje}
Archivo: \code{src/java/model/SolicitudMensaje.java}.
Responsabilidad: modelar mensajes/comentarios asociados a una solicitud (emisor, texto, adjunto, fecha) para mantener histórico de comunicaciones.

\subsection*{TipoSolicitud}
\addcontentsline{toc}{subsection}{TipoSolicitud}
Archivo: \code{src/java/model/TipoSolicitud.java}.
Responsabilidad: representar las categorías de trámites disponibles (nombre, descripción, responsable asignado).

\subsection*{ProgramaAcademico, Sede, Jornada}
\addcontentsline{toc}{subsection}{ProgramaAcademico / Sede / Jornada}
Archivos: \code{ProgramaAcademico.java}, \code{Sede.java}, \code{Jornada.java}.
Responsabilidad: tablas maestras que normales los valores usados en formularios y filtros.

% =================================================
\section{Patrón MVC y Flujo de Interacciones}
% =================================================

\subsection{Implementación del patrón MVC}

El proyecto aplica el patrón Modelo-Vista-Controlador (MVC) de forma clásica:

\begin{itemize}
    \item Modelos: clases en \code{src/java/model/} y DAOs en \code{src/java/dao/} que representan y persisten el estado de la aplicación.
    \item Vistas: archivos JSP ubicados en \code{web/} que reciben atributos desde los controladores y renderizan HTML.
    \item Controladores: servlets en \code{src/java/controller/} que procesan peticiones HTTP, validan entrada, invocan DAOs y despachan a JSP.
\end{itemize}

\subsection{Secuencia típica: creación de una solicitud (request)}

\begin{enumerate}
    \item El usuario en el navegador completa el formulario y hace POST a \code{/student/new-request}.
    \item \code{NewRequestServlet} recoge parámetros, valida datos y, si incluye archivo adjunto, llama a utilitarios para almacenar el fichero en \code{/uploads}.
    \item El servlet invoca al DAO responsable (por ejemplo, \code{SolicitudDAO} o métodos en \code{StudentDAO}) para insertar el registro en la BD y obtener el \code{solicitudId}.
    \item El servlet redirige al detalle de la solicitud o al listado del estudiante (\code{/student/request-detail?id=<solicitudId>} o \code{/student/requests}).
    \item La vista JSP correspondiente lee atributos de request/session y muestra la información actualizada.
\end{enumerate}

\subsection{Interacción Admin-Student (respuesta a solicitud)}

\begin{enumerate}
    \item El administrador abre el listado en \code{/admin/requests} y selecciona una solicitud.
    \item \code{AdminRequestDetailServlet} carga la solicitud y su historial (vía \code{SolicitudMensajeDAO}) y presenta la vista de detalle.
    \item El administrador edita el estado o escribe una respuesta; al enviar, el servlet guarda la respuesta en \code{solicitud\_mensaje} y actualiza el estado en \code{solicitud}.
    \item Si hay adjuntos en la respuesta, el servlet valida la extensión y tamaño, y los guarda en \code{/uploads/messages}.
    \item Opcionalmente, el código puede notificar por correo al estudiante (la implementación por defecto imprime en consola).
\end{enumerate}

\subsection{Filtros y seguridad}

El proyecto usa filtros HTTP en \code{src/java/filter/} para:
\begin{itemize}
    \item Verificar sesión y permisos antes de acceder a rutas protegidas.
    \item Prevenir CSRF básico (si aplica) y sanitizar entradas comunes.
\end{itemize}

\subsection{Manejo de archivos y restricciones}

\begin{itemize}
    \item Los servlets con \code{@MultipartConfig} definen \code{maxFileSize} y \code{maxRequestSize} (ver implementaciones: \code{NewRequestServlet}, \code{EditRequestServlet}, \code{RequestDetailServlet}, \code{AdminRequestDetailServlet}).
    \item Los controladores validan el tipo MIME y la extensión (p. ej., PDF para estudiantes, PDF/DOCX/JPG/PNG para admin).
\end{itemize}

\subsection{Buena práctica para contribuciones}

\begin{itemize}
    \item Añadir tests unitarios para lógica en utilidades y DAOs si se extrapola a JUnit.
    \item Mantener consultas SQL en el DAO y evitar lógica de negocio en JSP.
    \item Documentar cualquier cambio en \code{sql/database.sql} y actualizar el diagrama ER.
\end{itemize}

% =================================================
\section{Resumen de Interacciones entre Componentes}
% =================================================

\begin{longtable}{|p{0.28\textwidth}|p{0.65\textwidth}|}
\hline
Componente & Resumen de interacción \\
\hline
Servlets (controladores) & Reciben la petición HTTP, validan, llaman a DAOs/utiles y despachan a JSP. \\
\hline
DAOs & Ejecutan consultas SQL, mapean ResultSet a objetos de modelo y devuelven DTOs a servlets. \\
\hline
Models & Objetos-PONOs que transportan datos entre DAO y vista/servlet. \\
\hline
JSPs (vistas) & Muestran datos y formularios; deben evitar lógica de negocio compleja. \\
\hline
Filters & Interceptan peticiones para seguridad, logging y control de acceso. \\
\hline
Uploads folder & Ubicación física relativa dentro de la aplicación donde se guardan adjuntos; los servlets devuelven rutas relativas para su descarga. \\
\hline
\end{longtable}

\vspace{0.5cm}
\subsection*{Lecturas recomendadas para desarrolladores}
\addcontentsline{toc}{subsection}{Lecturas recomendadas}
\begin{itemize}
    \item Java Servlet Specification (documentación oficial).
    \item JDBC and Connection Pooling best practices.
    \item OWASP guidelines para manejo de archivos y control de accesos.
\end{itemize}

\end{document}
